%%%%%%%%%%%%%%%%%%%%%%%%%%%%%%%%%%%%%%%%%%%%%%%%%%%%%%%%%%%%%%%%%%%%%%%%%%%%%%
% Este é o modelo dun artigo simplificado. NON ESTÁ COMPLETO. ESTO NO É A
% REVISTA FINAL. A revista final componse de varios arquivos que inclúen o que
% aquí é o corpo do documento (dentro do \begin{document}). Esto so sirve para
% poder compilar un artigo individual facilmente, pero non inclue a portada,
% contraportada, e varios macros non están definidos. Esto no meu sistema
% compílase nuns 10 s. A revista_001 completa, en case 30 s
%
% Para usalo fai falta poñer 3 cousas no mesmo diretorio
% a clase da revista:   revista.cls
% este artigo:          modelo_artigo.tex
% as tipografias:       fontes\
%
% Opcionalmente, tamén é util ter o arquivo de configuracion 'latexmkrc', pero
% non é necesario.
%
% COMPILACION:
% Necesítase LuaLaTeX, non funciona doutro modo (nin o imos cambiar). Tamén se
% precisa Biber para a bibliografía
%
% Na maioría de editores debería bastar con premer calquer botón de 'executar'.
% Se non funciona seguramente teñades que cambiar o método que usa voso editor
% detrás de cámaras. Recordo que TeXStudio tiña un panel que deixaba
% seleccionar qué compiladores/executables queremos usar. O mellor é usar
% 'latexmk', o cal vai, á sua vez, elixir qué compiladores/executables se usan
% de xeito automático
%%%%%%%%%%%%%%%%%%%%%%%%%%%%%%%%%%%%%%%%%%%%%%%%%%%%%%%%%%%%%%%%%%%%%%%%%%%%%%

% Esto importa a clase definida no arquivo aparte, 'revista.cls'

% Este macro hai que usalo sempre. Fai bastantes cousas, como engadir o titulo
% e autor ao índice, resetear os conteos de figuras, e evidentemente imprimir
% no PDF a propia información de cada artigo.
%
% \Titular   -> sen asterisco, non engade a sección ao indice
% \Titular*  -> con asterisco, engade a sección ao indice.
%
% Ao ser unha versión simplificada da revista, da igual cal usedes. Eso xa se
% cambiará logo na propia edición
%
% {TITULO}   Evidente, o titulo do artigo. Non debería ser moi longo.
% {PERSONA}  Autores do artigo
% {TEMA}     A escoller entre:
%               divulgacion,historia,actualidadeFacultade,
%               actualidadeCientifica,filosofia,entrevistas,
%               programacion,pasatempos,anuncios,profesorado
% {SUBTITULO} Un pequeno texto algo máis explicativo (e posiblemente longo) co
%             titulo. E o único argumento do titular que pode quedar en branco
\Titular%
{José Edelstein: <<Quen facemos ciencia somos humanos, cada un cos seus erros e acertos>>}%
{Luis Arcas Morcillo}%
{entrevistas}%
{Por mor do centenario da visita de Einstein á Arxentina, falamos cun dos homes que máis sabe do pai da relatividade sobre á importancia da memoria na ciencia, a situación internacional da investigación e o panorama da divulgación.}%


% 'refsection' sirve para que ao final, ao mostrar a bibliografía, esta solo
% inclua as entradas que foron citadas DENTRO desta 'refsection' con \cite. É
% opcional, a menos que usades bibliografía, entón é obligatorio
\begin{refsection}

% O CORPO DO ARTIGO COMEZA DE AQUÍ EN DIANTE
%%%%%%%%%%%%%%%%%%%%%%%%%%%%%%%%%%%%%%%%%%%%%%%%%%%%%%%%%%%%%%%%%%%%%%%%%%%%%%
%vvvvvvvvvvvvvvvvvvvvvvvvvvvvvvvvvvvvvvvvvvvvvvvvvvvvvvvvvvvvvvvvvvvvvvvvvvvvv

% 'multicols' danos bastante flexibilidade para colocar as cousas en duas
% columnas. Se unha parte non debe ter 2 columnas, solo hai que rematar o
% entorno 'multicols' antes, escribir o que queiramos (qedará en 1 columna), e
% volver iniciar outro entorno 'multicols' para o resto
\begin{multicols}{2}


% Os únicos "niveis" que se permiten son de 'subsection' en abaixo. En
% principio poden estar numerados ou no, a gusto de cada un


Entramos no IGFAE e recíbenos cun sorriso na boca, aínda ten que acabar papelame (no momento que lle entrevistamos, recoñece que enviou as notas de Métodos Matemáticos VI ás 3:30 horas), e pídenos esperar un minuto. José, coñecido pola maioría do alumnado da Facultade, é un home atarefado, mais aínda ten tempo para atendernos e responder amablemente as nosas preguntas.\\

\begin{figure}[H]
    \centering
    \includegraphics[width=\linewidth]{revistas/004/imaxes/edelstein1.jpg}
    \caption{José Edelstein}
    \label{fig.: edelstein}
\end{figure}

Nacido en Buenos Aires, a súa formación académica e posterior experiencia no ámbito investigador e divulgativo fai conta dunha carreira internacional impresionante. No entanto, lémbranos que non sempre estivo vinculado á Física. \textbf{José comezou a licenciatura en Enxería Eléctrica en Venezuela}, a causa do exilio durante a ditadura do PRN na Arxentina. Volveu ó seu país natal en 1985 para proseguir cos seus estudos, pero a piques de rematar (no cuarto curso de cinco), deuse conta de que a súa paixón era a Física, polo que \textbf{deixou a licenciatura e foi aceptado no Instituto Balseiro}. Dende alí, unha apaixonante carreira académica arrancou: doutorando na Universidad de La Laguna, \textit{post-docs} en Santiago, Lisboa e Harvard; así como labores de investigación no CONICET (análogo ó CSIC na Arxentina) e no \emph{Centro de Estudios Científicos} (CEC) en Chile. Ata que finalmente a concesión da beca \emph{Ramón y Cajal} concédelle aquí, en Santiago, a estabilidade que todo investigador (e pai de dous nenos) busca.\\

Leva xa \textbf{máis de vinte anos vinculado á USC e ó IGFAE}, pero non esquece as súas raíces, colaborando en actividades culturais e divulgativas na Arxentina e outros países hispanoamericanos (exemplo notable é a iniciativa \emph{Universo entre Canciones}, cuxa xira internacional tivera parada aquí en Santiago hai uns anos). Un dos máis recentes fora a \textbf{celebración do centenario da visita de A. Einstein á Arxentina e ó Uruguai} nunha xira internacional por mor do seu Nobel, evento no cal Edelstein participou activamente como orador e organizador.

\begin{center}
    ***
\end{center}

\noindent \textbf{PREGUNTA.} \hspace{5px} \textbf{Por que é importante seguir celebrando as visitas de Einstein no mundo?}\\

\noindent \textbf{RESPOSTA.} \hspace{5px} Eu creo que a palabra clave é \textbf{<<legado>>}. Primeiro está o legado que deixa ós físicos: non hai día no que non apareza o seu nome nalgún \emph{paper} ou texto que consultes. Tamén está o legado que nós deixamos de Einstein: se non celebrásemos o día de Einstein coa importancia que merece e non nos limitásemos a simples efemérides en redes sociais, transmitimos a mensaxe de que quizais non é unha figura tan relevante. Nada máis lonxe da realidade: \textbf{Einstein foi un científico moi comprometido cultural e politicamente}. Foi pioneiro nos debates de desarmamento nuclear e participou activamente en discusións polos dereitos civís. Ademais, \textbf{eu conecto profundamente con el polo vínculo que nos une}: somos ambos os dous xudeus ateos, os meus avós tiveran que emigrar na Segunda Guerra Mundial para salvar a vida, ó igual que fixo Einstein. Así que espero que dentro de cen anos outro louco coma min se anime a organizar o segundo centenario.\\

\noindent \textbf{P.} \hspace{5px} \textbf{Cres que é necesario facer máis exercicios de memoria histórica na ciencia?} Isto é, poñer en valor figuras inxustamente esquecidas ou descubrimentos/avances relevantes que pasan desapercibidos normalmente.\\

\noindent \textbf{R.} \hspace{5px} Nós, como científicos, temos unha mentalidade distinta ós literatos e humanistas: \textbf{centrámonos máis na obra que no autor}. %Por exemplo, Einstein equivocouse moitas veces, mais ninguén lle segue soamente pola súa popularidade. Xa no seu momento, moitos científicos como Bohr, Pauli, aínda que non tivesen a reputación da que agora gozan, non dubidaron en argumentar con Einstein en torno á mecánica cuántica.
Ocórreseme o seguinte exemplo: Einstein propuxo a teoría da relatividade xeral, que outros autores adaptaron e escribiron cunha linguaxe máis limpa e precisa. Cando un estudante novo estuda relatividade xeral, baséase nos textos que outros autores redactaron en base ás ideas e traballos de Einstein: se vas directamente ós \textit{papers} orixinais, a notación é máis confusa e seguramente haberá erros que máis tarde serían corrixidos. Isto pasa menos nas humanidades. Quen mellor expresa o ideario de Heidegger é o propio Heidegger, polo que á hora de estudar a súa filosofía un recorre directamente ós libros do autor. Moi posiblemente é por iso que se tende a venerar máis o autor que a obra en si, cousa que en ciencias naturais é menos común.

Non obstante, como humanos, non podemos esquecer que \textbf{quen facemos ciencia tamén somos humanos}, cada un coa súa biografía, coas súas dúbidas, coa súa tenacidade e cos seus erros e acertos. Facer exercicios de memoria na ciencia axuda a poñer en valor esta faceta e a humanizar a ciencia, o cal creo que é importante.\\

\noindent \textbf{P.} \hspace{5px} Na \emph{laudatio} que pronunciaches na charla de Kip Thorne no contexto do programa \textit{ConCiencia} nomeaches un quinteto de científicos que son clave para entender a relatividade xeral: \textbf{Minkowski, Einstein, Penrose, Thorne e Hawking}. Invítoche a facer un exercicio de memoria histórica: \textbf{hai algún científico máis que engadirías a esta lista} de grandes científicos, non só no contexto da relativade xeral, senón dentro doutros ámbitos; \textbf{que tivese marcado a túa traxectoria?}\\

\noindent \textbf{R.} \hspace{5px} Sen ir máis lonxe, o posto de Thorne podería intercambiarse con \textbf{John Wheeler}, o seu director de tese e quen lle introduciu na relatividade xeral. El puxera a brillantes alumnos como Feynman, Thorne ou Misner a traballar nun campo que nos anos 60 non era tan relevante, permitindo un grande avance no campo ata os nosos días. Thorne díxome que tería posto o seu mestre no seu lugar, pero finalmente considero que Thorne vai alí.

Outro científico que vén á miña mente é \textbf{Paul Dirac}. Para min, o segundo científico máis brillante do século XX: unha figura moi creativa e cunha biografía moi interesante. Tamén podería mencionar a \textbf{Richard Feynman, Wolfgang Pauli...} Como divulgador, asomarse á vida destes xenios revela outros personaxes, quizais menores ó lado destes, pero moi interesantes. Como científico, un pode coñecer mentes brillantes actuais, tocadas por unha variña máxica. Por exemplo, \textbf{Juan Martín Maldacena}: compañeiro e amigo na licenciatura e gran figura da física teórica mundial, con grandes avances na formulación cuántica da gravidade. E xa que estamos na USC, quería reivindicar a figura de \textbf{Marcos Mariño Beiras}. Actualmente é profesor de Física Matemática na Universidade de Xenebra. Home de grande intelixencia matemática e intuitiva á hora de formular hipóteses. Aí tedes idea para unha futura entrevista.\\

\noindent \textbf{P.} \hspace{5px} Dende logo que si, grazas pola suxestión \emph{(Risas)}. Seguindo coa entrevista, quería agora preguntarche polo programa \emph{ConCiencia}, dirixido por Jorge Mira. Nel, ademais dos vínculos que tende o encontro entre o equipo do convidado e a universidade, hai unha charla pública na que o relator fala dun tema concreto. Ás veces, pode resultar complicado seguir o tema, especialmente para a audiencia non especializada. \textbf{Que recomendarías a quen asiste a unha destas charlas e non comprende o tema ou se perde durante ela?} \\

\noindent \textbf{R.} \hspace{5px} Eu son da opinión de que \textbf{as charlas teñen que ser desafiantes cos asistentes}. Hai distintas escolas na divulgación, e cada cal escolle as súas preferencias á hora de expoñer un tema. En calquera caso, unha calidade indispensable é a \textbf{claridade}, e por exemplo, Kip [na charla da última edición de ConCiencia], foi moi claro, aínda que [a charla] fose un pouco máis técnica. Neste sentido, un ten feito un bo traballo se quen escoita a charla marcha con algunha dúbida e, máis importante, con gana de encontrar a resposta. \textbf{Eu recomendaría que a xente tomase nota de posibles referencias}: películas relacionadas co tema, libros, outros relatorios,... Se vas dar unha charla sinxela, onde todos saben de todo o que se vai falar, entón nunca traes un nobel.\\

\noindent \textbf{P.} \hspace{5px} Recapitulo un momento nos proxectos divulgativos nos que estás involucrado: \emph{Amautas}, unha plataforma dixital de aprendizaxe; \emph{Coffee Break}, un podcast de actualidade científica; \emph{Universo entre Canciones}, un recital onde música e ciencia se fan unha. No relatorio de Kip Thorne, el dixera que a arte era unha das súas paixóns, e que gustaba de usalas para transmitir a ciencia. \textbf{Ti cres que a chave da divulgación está hoxe na arte} en xeral: pictórica, musical, teatral, dixital...?\\

\noindent \textbf{R.} \hspace{5px} Si, nunha charla que dou recorrentemente chamada \emph{Vasos comunicantes} reflexiono sobre o vehículo que supón a arte para comunicar a ciencia, e a inspiración que pode ser a ciencia para a arte. Ademais, \textbf{existen paralelismos entre ciencia e arte}: cando un pintor está a comezar un cadro, non pensa, <<Que dirá o público da miña obra?>>; senón que primeiro o remata e despois xa recibe as críticas. Na ciencia pasa o mesmo. No entanto, a simbiose ciencia-arte non é sempre tan obvia: hai actividades que eu organizo en espazos públicos (teatros, auditorios...) que a xente pregunta se está reservado ó público especializado. Ninguén preguntaría iso a un artista. \textbf{A ciencia}, aínda que máis técnica e críptica, \textbf{tamén é para o público, non deixa de ser unha actividade humana}. Hoxe temos unha gran diversificación da oportunidade: se pensas en calquera medio que che permita divulgar un tema, seguramente xa haxa alguén que estea a usalo.

\begin{figure}[H]
    \centering
    \includegraphics[width=\linewidth]{revistas/004/imaxes/universo_canciones.jpg}
    \caption{José (centro) e Lore Edelstein (dereita), xunto a Daniela de Rito (esquerda), integrantes de \emph{Universo entre Canciones} (Fonte: José Edelstein).}
    \label{fig.: universo_canciones}
\end{figure}

En relación á divulgación en redes, quería deixar un recado á xente máis nova que queira dedicarse á divulgación: \textbf{centrádevos en rematar a titulación e en facer dous ou tres post-docs}. É entón cando podedes comezar a divulgar activamente. Cando un se lanza a explicar un tema que aprendeu o día anterior, todo está collido con alfinetes; e se alguén sopra máis ou menos forte, todo se afunde. \textbf{É necesario ter coñecementos profundos da materia}, e lixeiros cambios na explicación, por moi sutís que poidan parecer, poden cambiar totalmente o sentido conceptual do que estás a explicar. Non esquezamos que os grandes divulgadores actuais e de sempre son persoas cunha grande experiencia académica. É como Messi: fai pensar que é moi sinxelo xutar á portería e marcar gol, pero detrás hai horas e horas de adestramento. O talento innato é diferencial, dende logo, pero só con el non vas chegar moi lonxe.\\

\noindent \textbf{P.} \hspace{5px} Quería pararme agora en \emph{Amautas}, onde polo pago dunha subscrición tes acceso a cursos impartidos por diversos profesores e figuras referentes na divulgación actual: Santaolalla e ti mesmo (cofundadores da plataforma), Jorge Mira, Lemnismath... Notas que o interese pola ciencia en xeral e a física en particular crecera nos últimos anos? \textbf{Anímase máis a xente a estudar física dentro e fora das aulas?}\\

\noindent \textbf{R.} \hspace{5px} A nivel académico é notable o aumento da nota de corte nos últimos 15 anos. \textbf{Pasamos de non ter nota de corte a chegar a case 12 sobre 14 aquí en Santiago}. Pasou de haber promocións nas que todos entraron con Física como a súa segunda ou terceira opción a achar unha demanda moi superior á oferta (este ano, as solicitudes de admisión ó grao alcanzaron as catrocentas aproximadamente). Este interese pola física é grazas a figuras públicas na divulgación como por exemplo Santaolalla: o seu público máis novo proxéctase neles e busca emular os seus referentes.

Non obstante, permíteme ser cauto na miña resposta, porque tamén observo un aumento do escurantismo por parte de certos sectores, especialmente despois da pandemia da COVID-19. \textbf{Fai falta conseguir que a cidadanía sexa crítica} grazas a unha formación en conceptos básicos da ciencia a fin de evitar malentendidos e desinformación á hora de interpretar resultados de enquisas, estudos, etc. Por exemplo, é moi importante que a xente entenda que \textbf{a ciencia se basea en preguntas}, e que \textbf{a verdade que manexamos non é absoluta}, senón que pode cambiar parcial ou totalmente, e debe cambiar se os resultados o piden. E non por iso a ciencia perde credibilidade, senón que é máis obxectiva.\\

\noindent \textbf{P.} \hspace{5px} Dende logo, o espírito crítico é fundamental na sociedade, máis a día de hoxe. Con toda a túa actividade divulgadora, \textbf{como compaxinas a túa carreira investigadora e docente na USC?}\\

\noindent \textbf{R.} \hspace{5px} É moi difícil, quitando horas ó sono e á familia. Aínda que tivera incursións moi puntuais previamente, eu non empecei o meu labor divulgativo ata ser profesor na universidade. A miña experiencia á hora de escribir e falar en conferencias deu pé a un aumento de tempo investido nesta actividade. Ademais, fronte ó panorama actual, onde figuras disruptivas están a negar a importancia da ciencia na sociedade en moitos países (e lamentablemente na miña Arxentina tamén ocorre); preséntase ante min unha conxuntura: debo centrarme máis na investigación en temas vangardistas e moi restrinxidos, ou me dedico a facer chegar información máis xeral pero moi relevante á cidadanía? Eu considérome unha persoa politicamente comprometida, e \textbf{comunicar ciencia é para min e un compromiso coa sociedade na que vivo}. Isto implica moito sacrificio, e estes últimos anos a miña actividade investigadora reduciuse en favor da miña carreira divulgadora. Sen ir máis lonxe, este ano espero publicar dous libros; pero espero recuperar pronto o \emph{momentum} na investigación.

No meu parón investigador tamén puiden <<abrir o espectro>>: para preparar o programa \emph{Coffee Break}. Teño que estar informado das novidades en distintas ramas da ciencia, e descubrín no proceso enfoques distintos a problemas que me interesan do meu campo de investigación, e que espero poder aplicar nun futuro. Esta retroalimentación motívame a seguir divulgando.\\

\noindent \textbf{P.} \hspace{5px} Parémonos nun instante nas <<figuras disruptivas>> das que falabas antes. Claramente sabemos a quen nos estamos a referir, e non é unha opinión illada. Lembro que Kip Thorne na súa charla en Santiago afirmou que nos Estados Unidos a educación sofre, e que é o momento de Europa de converterse na defensora do progreso científico. Na Arxentina a situación tampouco é mellor. Lamentablemente, España comparte coa Arxentina unha fuga de cerebros que parece complicada de paliar; máis estes últimos anos. \textbf{Que virtudes ten o sistema español que se poderían aplicar na Arxentina e viceversa?} \\

\noindent \textbf{R.} \hspace{5px} Unha pregunta difícil, dar contexto á situación na Arxentina levaríanos a unha entrevista de cinco horas como mínimo. España ten unha estabilidade institucional, monetaria, etc. que non hai na Arxentina. Isto obriga a gran maioría de colegas físicos a emigrar, mentres que aquí en España non me atrevería a afirmar que a maioría marchan ó estranxeiro. Non obstante, como afirma un dos meus discípulos, que está a traballar no Canadá, o sistema español ten precariedade nalgúns aspectos.

\textbf{Dende España estanse a promover diferentes medidas para fomentar a estabilidade do talento nacional}: bolsas \emph{Ramón y Cajal}, institucións como \emph{Icrea} en Barcelona ou \emph{Ikerbasque} no País Vasco... Búscase desenvolver programas en Galicia similares ó \emph{Icrea}. Moitas opcións, moi competitivas, certo; pero hai onde elixir.

Non obstante, na Arxentina a única vía é aplicar a carreira investigadora polo CONICET, e descoñezo se segue a haber convocatorias. Aínda se accedes, os salarios non che permiten dedicarte exclusivamente á ciencia. Á hora de traballar nun laboratorio e pedir orzamentos para experimentos, \textbf{en España o diñeiro chega en tempo e forma} - un euro vale considerablemente o mesmo dun ano para outro. \textbf{Na Arxentina}, o valor do peso cambia da noite ó día, e normalmente o Estado demórase nos pagos, polo que \textbf{moitas veces o orzamento inicial resulta insuficiente}. Dado que tes que entregar o informe final igualmente, os poucos valentes que quedan teñen que <<atar todo con arame>>, conseguindo montar un laboratorio de 5000€ con só 5€, por exemplo.\\

\noindent \textbf{P.} \hspace{5px} Grazas pola túa reflexión, José. A situación en ambos os dous países é complicada, pero dende logo aquí en España temos sorte de habitar no contexto europeo, onde o valor de institucións como o CERN son exemplo de que a ciencia se valora por enriba de ideais, un oasis entre os desertos da tensión xeopolítica.\\

\noindent \textbf{R.} \hspace{5px} \textbf{Os europeos deberían saber a sorte que teñen co CERN. Eu daríalle á institución o Nobel da Paz}, soamente por conseguir que xente de diferentes nacionalidades coopere polo avance do coñecemento en beneficio da humanidade.\\

\noindent \textbf{P.} \hspace{5px} Xa para ir rematando a entrevista, quería lanzar unha pregunta atrevida: ti fuches patrón na visita de Stephen Hawking a Santiago en 2008. \textbf{Cres que nalgún momento a figura de Hawking será tan relevante coma a de Einstein?} Celebraránse as súas visitas polo mundo igual que facemos con el?\\

\noindent \textbf{R.} \hspace{5px} Eu creo que non, Hawking estaría de acordo comigo que Einstein é moito máis importante. Obviando o seu legado, a súa figura, ademais de excéntrica, está cuberta por un contexto histórico singular (guerras mundiais, nacemento da mecánica cuántica...). \textbf{Hawking está no panteón dos grandes físicos} polo seu extenso traballo: aclarou moitas dúbidas sobre a orixe do universo e a súa estrutura, así como dos buratos negros. Non obstante, coetáneos que traballaron nos mesmos campos, como Penrose, son menos coñecidos. O salto á cultura pop vén da súa discapacidade: a súa grande intelixencia a pesar dunha enfermidade terminal inspirou a súa aparición en series como \textit{Os Simpsons}, \textit{The Big Bang Theory}, \textit{Star Trek}... \textbf{Tamén é un exemplo de tenacidade e actitude}. A súa vontade de querer gozar a vida cada segundo pese ás adversidades debería inspirar a xente que se afunde fronte a problemas que son <<puro \emph{gauge}>>.

A figura de Hawking será lembrada, dende logo, pero dubido que en 2108 alguén faga un aniversario da súa visita a Santiago. Eu non estarei aquí para comprobalo, pero ti quizais si.\\

\noindent \textbf{P.} \hspace{5px} Vai ser complicado, non é doado chegar ós 105 anos \emph{(Risas)}. Pero escrito queda, por se alguén le esta entrevista dentro de cen anos.

% \begin{center}
%     ***
% \end{center}

% Novamente agradecemos a José o seu tempo e amabilidade antes de saír pola porta. O mundo é imperfecto, tamén a ciencia (é o que ten ser feita por humanos). Pero grazas ó espírito crítico infundido por persoas como Edelstein, un pode afrontar o futuro con tenacidade fronte a adversidade, ó menos ata o próximo centenario da visita de Einstein.



% Finalmente mostrar a bibliografía. Recordo que é importantísimo que esto esté
% dentro do entorno 'refsection'. Engadir que a bibliograia manéxase con
% Biblatex
%\printbibliography

\end{multicols}
\end{refsection}
